\documentclass[12pt,a4paper]{article}
\usepackage[utf8]{inputenc}
\usepackage[T1]{fontenc}
\usepackage[brazil]{babel}
\usepackage{amsmath,amssymb,amsfonts}
\usepackage{siunitx}
\usepackage{geometry}
\usepackage{graphicx}
\usepackage{float}
\usepackage{xcolor}
\usepackage{hyperref}
\usepackage{listings}
\usepackage{booktabs}
\usepackage{enumitem}
\geometry{left=3cm,right=2cm,top=3cm,bottom=2cm}
\hypersetup{colorlinks=true, linkcolor=black, urlcolor=blue, citecolor=black}
\lstdefinestyle{pythonstyle}{
    language=Python,
    basicstyle=\ttfamily\footnotesize,
    keywordstyle=\color{blue},
    commentstyle=\color{green!40!black},
    stringstyle=\color{red!60!black},
    numbers=left,
    numberstyle=\tiny,
    stepnumber=1,
    numbersep=8pt,
    showstringspaces=false,
    breaklines=true,
    frame=single,
    tabsize=4
}
\title{Solução -- Primeira Avaliação de Métodos Numéricos}
\author{Gustavo Bittar Gonçalves}
\date{}

\begin{document}
\maketitle

\section*{Observação inicial}
Este documento apresenta uma solução-modelo para a avaliação de Métodos Numéricos, com ênfase na formulação matemática, discretização por diferenças finitas e implementação computacional em Python. Onde o enunciado não fornece explicitamente propriedades termofísicas do material, adotam-se valores usuais para alumínio puro, declarados na seção numérica.

\section{Questão 1 -- Equação governante para $k$ constante}
Para condução térmica bidimensional, transiente, sem geração interna de calor e com propriedades constantes, a equação diferencial parcial governante é:
\begin{equation}
\rho c_p \frac{\partial T}{\partial t} = k\left(\frac{\partial^2 T}{\partial x^2}+\frac{\partial^2 T}{\partial y^2}\right),
\qquad 0<x<L_x,\;0<y<L_y,\;t>0.
\end{equation}

Também pode ser escrita como:
\begin{equation}
\frac{\partial T}{\partial t} = \alpha \left(\frac{\partial^2 T}{\partial x^2}+\frac{\partial^2 T}{\partial y^2}\right),
\qquad \alpha = \frac{k}{\rho c_p}.
\end{equation}

As condições de contorno do enunciado ficam:
\begin{align}
T(0,y,t) &= T_{\text{imp}}(y,t), && 0<y<L_y, \\
\left.\frac{\partial T}{\partial x}\right|_{x=L_x} &= 0, && 0<y<L_y, \\
-k\left.\frac{\partial T}{\partial y}\right|_{y=0} &= h_1\big(T(x,0,t)-T_{\infty,1}\big), && 0<x<L_x, \\
-k\left.\frac{\partial T}{\partial y}\right|_{y=L_y} &= h_2\big(T(x,L_y,t)-T_{\infty,2}\big) + \varepsilon\sigma\left(T(x,L_y,t)^4-T_{\text{viz}}^4\right), && 0<x<L_x.
\end{align}

A condição inicial é:
\begin{equation}
T(x,y,0)=300\,\si{K}, \qquad 0<x<L_x,\;0<y<L_y.
\end{equation}

\section{Questão 2 -- Equação governante para $k=aT+b$}
Se a condutividade térmica varia linearmente com a temperatura,
\begin{equation}
k(T)=aT+b,
\end{equation}
a forma conservativa correta da equação governante é:
\begin{equation}
\rho c_p \frac{\partial T}{\partial t} = \frac{\partial}{\partial x}\left(k(T)\frac{\partial T}{\partial x}\right) + \frac{\partial}{\partial y}\left(k(T)\frac{\partial T}{\partial y}\right).
\end{equation}

Substituindo $k(T)=aT+b$:
\begin{equation}
\rho c_p \frac{\partial T}{\partial t} = \frac{\partial}{\partial x}\left[(aT+b)\frac{\partial T}{\partial x}\right] + \frac{\partial}{\partial y}\left[(aT+b)\frac{\partial T}{\partial y}\right].
\end{equation}

As condições de contorno e a condição inicial permanecem com a mesma estrutura da Questão 1, mudando apenas o valor local de $k$ quando necessário.

\section{Questão 3 -- Diferenças finitas de segunda ordem no espaço e Euler no tempo}
Aplicando o método de Euler explícito no tempo e diferenças finitas centrais de segunda ordem no espaço para a equação da Questão 1, obtém-se:
\begin{equation}
\frac{T_{i,j}^{n+1}-T_{i,j}^{n}}{\Delta t} = \alpha \left[
\frac{T_{i+1,j}^{n}-2T_{i,j}^{n}+T_{i-1,j}^{n}}{\Delta x^2}
+
\frac{T_{i,j+1}^{n}-2T_{i,j}^{n}+T_{i,j-1}^{n}}{\Delta y^2}
\right].
\end{equation}

Isolando $T_{i,j}^{n+1}$:
\begin{equation}
T_{i,j}^{n+1}=T_{i,j}^{n}
+ \alpha\Delta t \left[
\frac{T_{i+1,j}^{n}-2T_{i,j}^{n}+T_{i-1,j}^{n}}{\Delta x^2}
+
\frac{T_{i,j+1}^{n}-2T_{i,j}^{n}+T_{i,j-1}^{n}}{\Delta y^2}
\right].
\end{equation}

Definindo:
\begin{equation}
Fo_x = \frac{\alpha\Delta t}{\Delta x^2},
\qquad
Fo_y = \frac{\alpha\Delta t}{\Delta y^2},
\end{equation}
a equação de atualização pode ser escrita como:
\begin{equation}
T_{i,j}^{n+1}=(1-2Fo_x-2Fo_y)T_{i,j}^{n} + Fo_x(T_{i+1,j}^{n}+T_{i-1,j}^{n}) + Fo_y(T_{i,j+1}^{n}+T_{i,j-1}^{n}).
\end{equation}

Para uma malha uniforme com $\Delta x=\Delta y$, a condição usual de estabilidade do esquema explícito é:
\begin{equation}
Fo_x + Fo_y \leq \frac{1}{2}.
\end{equation}

\subsection*{Tratamento das condições de contorno}
\begin{itemize}[leftmargin=1.2cm]
    \item \textbf{Em $x=0$}: temperatura prescrita, logo $T_{0,j}^{n}=400\,\si{K}$ no caso numérico.
    \item \textbf{Em $x=L_x$}: adiabático, logo $\partial T/\partial x = 0$, o que pode ser aproximado por
    \begin{equation}
    \frac{T_{N_x,j}^{n}-T_{N_x-1,j}^{n}}{\Delta x}=0
    \quad\Rightarrow\quad
    T_{N_x,j}^{n}=T_{N_x-1,j}^{n}.
    \end{equation}
    \item \textbf{Em $y=0$}: convecção,
    \begin{equation}
    -k\left.\frac{\partial T}{\partial y}\right|_{y=0}=h_1(T-T_{\infty,1}).
    \end{equation}
    Usando diferença progressiva de primeira ordem na fronteira:
    \begin{equation}
    -k\frac{T_{i,1}^{n}-T_{i,0}^{n}}{\Delta y}=h_1(T_{i,0}^{n}-T_{\infty,1}),
    \end{equation}
    resultando em
    \begin{equation}
    T_{i,0}^{n}=\frac{\tfrac{k}{\Delta y}T_{i,1}^{n}+h_1T_{\infty,1}}{\tfrac{k}{\Delta y}+h_1}.
    \end{equation}
    \item \textbf{Em $y=L_y$}: convecção + radiação,
    \begin{equation}
    -k\left.\frac{\partial T}{\partial y}\right|_{y=L_y}=h_2(T-T_{\infty,2})+\varepsilon\sigma(T^4-T_{\text{viz}}^4).
    \end{equation}
    Como o termo radiativo é não linear, uma forma prática para o algoritmo explícito é usar a temperatura conhecida no tempo $n$ e definir um coeficiente radiativo equivalente:
    \begin{equation}
    h_{\text{rad}}^{n}=\varepsilon\sigma\left(T_{s}^{n}+T_{\text{viz}}\right)\left[(T_{s}^{n})^2+T_{\text{viz}}^2\right],
    \end{equation}
    de modo que
    \begin{equation}
    q_{\text{rad}}'' \approx h_{\text{rad}}^{n}(T_s^{n}-T_{\text{viz}}).
    \end{equation}
    Assim, a fronteira superior é tratada como uma convecção equivalente com
    \begin{equation}
    h_{\text{eq}}^{n}=h_2+h_{\text{rad}}^{n},
    \qquad
    T_{\infty,\text{eq}}^{n}=\frac{h_2T_{\infty,2}+h_{\text{rad}}^{n}T_{\text{viz}}}{h_2+h_{\text{rad}}^{n}}.
    \end{equation}
\end{itemize}

\section{Questão 4 -- Algoritmo}
Um algoritmo compatível com a discretização acima pode ser descrito da seguinte forma:

\begin{enumerate}[leftmargin=1.2cm]
    \item Definir $L_x$, $L_y$, $N_x$, $N_y$, $\Delta x$, $\Delta y$, $\Delta t$, $t_{\max}$ e as propriedades do material.
    \item Inicializar toda a malha com $T=300\,\si{K}$.
    \item Aplicar a condição de contorno prescrita em $x=0$.
    \item Para cada passo de tempo:
    \begin{enumerate}[label=\alph*)]
        \item atualizar a fronteira adiabática em $x=L_x$;
        \item atualizar a fronteira convectiva em $y=0$;
        \item atualizar a fronteira em $y=L_y$ com convecção + radiação linearizada;
        \item calcular as temperaturas internas com a equação explícita de Euler;
        \item sobrescrever os valores antigos pelos novos;
        \item armazenar campos e séries temporais nos instantes desejados.
    \end{enumerate}
    \item Ao final, gerar os gráficos pedidos no enunciado.
\end{enumerate}

\section{Questão 5 -- Solução numérica em Python}
Para a parte numérica, utiliza-se:
\begin{itemize}[leftmargin=1.2cm]
    \item $L_x=L_y=1.0\,\si{m}$;
    \item espessura $e=5.0\,\si{mm}$;
    \item $T(x=0,y,t)=400\,\si{K}$;
    \item em $x=L_x$: adiabático;
    \item em $y=0$: $h_1=100\,\si{W/(m^2.K)}$ e $T_{\infty,1}=350\,\si{K}$;
    \item em $y=L_y$: $h_2=200\,\si{W/(m^2.K)}$, $T_{\infty,2}=280\,\si{K}$, $\varepsilon=0.8$ e $T_{\text{viz}}=300\,\si{K}$;
    \item condição inicial: $T=300\,\si{K}$.
\end{itemize}

Como o enunciado não fornece explicitamente as propriedades termofísicas do alumínio puro, adota-se no código:
\begin{equation}
\rho = 2700\,\si{kg/m^3}, \qquad c_p = 900\,\si{J/(kg.K)}, \qquad k = 237\,\si{W/(m.K)}.
\end{equation}
Logo,
\begin{equation}
\alpha = \frac{k}{\rho c_p}.
\end{equation}

\subsection*{Escolha do passo de tempo}
Para o esquema explícito, usa-se um passo de tempo calculado automaticamente com fator de segurança, respeitando a estabilidade em função da malha.

\subsection*{Código Python comentado}
O código abaixo resolve a parte numérica, gera os gráficos pedidos e salva os arquivos PNG automaticamente.

\lstset{style=pythonstyle}
\begin{lstlisting}
import os
import numpy as np
import matplotlib.pyplot as plt

# ============================================================
# SOLUÇÃO NUMÉRICA DA CONDUÇÃO 2D TRANSIENTE EM UMA CHAPA
# Método: Euler explícito no tempo + diferenças finitas centrais
# Autor: Gustavo Bittar Gonçalves
# ============================================================

# ------------------------------
# Constantes físicas
# ------------------------------
sigma = 5.670374419e-8  # constante de Stefan-Boltzmann [W/m².K^4]

# ------------------------------
# Dados do problema
# ------------------------------
Lx = 1.0           # comprimento em x [m]
Ly = 1.0           # comprimento em y [m]
e = 5.0e-3         # espessura da chapa [m]

rho = 2700.0       # densidade do alumínio [kg/m³]
cp = 900.0         # calor específico [J/(kg.K)]
k = 237.0          # condutividade térmica [W/(m.K)]
alpha = k / (rho * cp)  # difusividade térmica [m²/s]

# Condições de contorno
T_left = 400.0     # temperatura prescrita em x=0 [K]
h_bottom = 100.0   # convecção em y=0 [W/(m².K)]
Tinf_bottom = 350.0

h_top = 200.0      # convecção em y=Ly [W/(m².K)]
Tinf_top = 280.0
epsilon = 0.8      # emissividade
Tviz = 300.0       # temperatura de vizinhança radiativa [K]

T_init = 300.0     # condição inicial [K]
t_final = 4.0 * 3600.0  # 4 horas em segundos

# Instantes pedidos para os campos de temperatura
field_times = [0.0, 1800.0, 3600.0, 7200.0, 10800.0, 14400.0]

# Malhas pedidas no enunciado para comparação
mesh_list = [10, 20, 40, 80]

# Pasta de saída para os PNGs
output_dir = "graficos_saida"
os.makedirs(output_dir, exist_ok=True)


def compute_stable_dt(dx, dy, alpha, safety=0.45):
    """
    Calcula um passo de tempo estável para o esquema explícito 2D.

    Para a discretização explícita em 2D:
        alpha*dt*(1/dx² + 1/dy²) <= 1/2

    Usamos um fator de segurança multiplicando o valor-limite.
    """
    dt_lim = 1.0 / (2.0 * alpha * (1.0 / dx**2 + 1.0 / dy**2))
    return safety * dt_lim


def radiative_h(Ts, Tviz, epsilon, sigma):
    """
    Calcula o coeficiente radiativo equivalente:
        h_rad = epsilon*sigma*(Ts + Tviz)*(Ts² + Tviz²)

    Essa expressão lineariza o termo radiativo:
        epsilon*sigma*(Ts^4 - Tviz^4) = h_rad*(Ts - Tviz)
    """
    return epsilon * sigma * (Ts + Tviz) * (Ts**2 + Tviz**2)


def apply_boundary_conditions(T, dx, dy):
    """
    Aplica as condições de contorno diretamente na matriz T.
    A matriz é indexada como T[j, i], onde:
      - i varia em x
      - j varia em y

    Convenção adotada:
      j = 0      -> y = 0 (borda inferior)
      j = Ny-1   -> y = Ly (borda superior)
      i = 0      -> x = 0 (borda esquerda)
      i = Nx-1   -> x = Lx (borda direita)
    """
    Ny, Nx = T.shape

    # 1) x = 0 -> temperatura prescrita
    T[:, 0] = T_left

    # 2) x = Lx -> adiabático => dT/dx = 0 => T[:, -1] = T[:, -2]
    T[:, -1] = T[:, -2]

    # 3) y = 0 -> convecção
    # -k*(dT/dy) = h*(T_s - T_inf)
    # Aproximação com diferença finita de 1ª ordem:
    # T[0, i] = ((k/dy)*T[1, i] + h*T_inf) / ((k/dy) + h)
    for i in range(1, Nx - 1):
        T[0, i] = ((k / dy) * T[1, i] + h_bottom * Tinf_bottom) / ((k / dy) + h_bottom)

    # 4) y = Ly -> convecção + radiação linearizada
    #    Usamos a temperatura atual da superfície para calcular h_rad
    for i in range(1, Nx - 1):
        Ts_old = T[-1, i]
        h_rad = radiative_h(Ts_old, Tviz, epsilon, sigma)
        h_eq = h_top + h_rad
        Tinf_eq = (h_top * Tinf_top + h_rad * Tviz) / h_eq
        T[-1, i] = ((k / dy) * T[-2, i] + h_eq * Tinf_eq) / ((k / dy) + h_eq)

    # Reaplica os cantos de forma simples e consistente
    T[0, 0] = T_left
    T[-1, 0] = T_left
    T[0, -1] = T[0, -2]
    T[-1, -1] = T[-1, -2]

    return T


def solve_case(Nx, Ny, save_prefix="malha"):
    """
    Resolve o problema para uma dada malha Nx x Ny.
    Retorna coordenadas, histórico e campo final.
    Além disso, salva gráficos dos campos e resultados solicitados.
    """
    dx = Lx / (Nx - 1)
    dy = Ly / (Ny - 1)
    x = np.linspace(0.0, Lx, Nx)
    y = np.linspace(0.0, Ly, Ny)

    dt = compute_stable_dt(dx, dy, alpha, safety=0.45)
    n_steps = int(np.ceil(t_final / dt))
    dt = t_final / n_steps  # ajusta para atingir exatamente t_final

    # matriz de temperatura inicial
    T = np.full((Ny, Nx), T_init, dtype=float)
    T = apply_boundary_conditions(T, dx, dy)

    # pontos para histórico temporal
    ix_mid = np.argmin(np.abs(x - Lx / 2.0))
    iy_mid = np.argmin(np.abs(y - Ly / 2.0))
    ix_q = np.argmin(np.abs(x - Lx / 4.0))
    iy_q = np.argmin(np.abs(y - Ly / 4.0))

    time_hist = [0.0]
    T_mid_hist = [T[iy_mid, ix_mid]]
    T_q_hist = [T[iy_q, ix_q]]

    # snapshots para os campos nos tempos pedidos
    snapshots = {0.0: T.copy()}

    # laço temporal principal
    for n in range(1, n_steps + 1):
        t = n * dt
        T_new = T.copy()

        # Atualização dos pontos internos via Euler explícito
        for j in range(1, Ny - 1):
            for i in range(1, Nx - 1):
                d2Tdx2 = (T[j, i + 1] - 2.0 * T[j, i] + T[j, i - 1]) / dx**2
                d2Tdy2 = (T[j + 1, i] - 2.0 * T[j, i] + T[j - 1, i]) / dy**2
                T_new[j, i] = T[j, i] + alpha * dt * (d2Tdx2 + d2Tdy2)

        # Aplica as condições de contorno após atualizar o interior
        T_new = apply_boundary_conditions(T_new, dx, dy)
        T = T_new

        # Guarda históricos temporais
        time_hist.append(t)
        T_mid_hist.append(T[iy_mid, ix_mid])
        T_q_hist.append(T[iy_q, ix_q])

        # Se t estiver suficientemente próximo de um tempo de interesse,
        # salva o campo de temperatura.
        for tf in field_times:
            if tf not in snapshots and abs(t - tf) <= 0.5 * dt:
                snapshots[tf] = T.copy()

    # Garantia de snapshot final
    snapshots[t_final] = T.copy()

    # Salva campos de temperatura como mapas de contorno
    X, Y = np.meshgrid(x, y)
    for tf in field_times:
        if tf in snapshots:
            plt.figure(figsize=(7, 5))
            cp_plot = plt.contourf(X, Y, snapshots[tf], levels=30)
            plt.colorbar(cp_plot, label="Temperatura [K]")
            plt.xlabel("x [m]")
            plt.ylabel("y [m]")
            plt.title(f"Campo de temperatura - t = {tf:.0f} s - malha {Nx}x{Ny}")
            plt.tight_layout()
            plt.savefig(os.path.join(output_dir, f"campo_T_t{int(tf)}_{save_prefix}_{Nx}x{Ny}.png"), dpi=300)
            plt.close()

    return {
        "x": x,
        "y": y,
        "T": T,
        "snapshots": snapshots,
        "time": np.array(time_hist),
        "T_mid_hist": np.array(T_mid_hist),
        "T_q_hist": np.array(T_q_hist),
    }


# ============================================================
# EXECUÇÃO PARA TODAS AS MALHAS PEDIDAS
# ============================================================
results = {}
for N in mesh_list:
    print(f"Resolvendo para malha {N}x{N}...")
    results[N] = solve_case(N, N, save_prefix="avaliacao1")

# ============================================================
# GRÁFICOS DE PERFIL EM t = 4 h
# ============================================================

def plot_vertical_profile(x_target, title_suffix, filename):
    plt.figure(figsize=(7, 5))
    for N in mesh_list:
        x = results[N]["x"]
        y = results[N]["y"]
        T = results[N]["T"]
        ix = np.argmin(np.abs(x - x_target))
        plt.plot(T[:, ix], y, label=f"{N}x{N}")
    plt.xlabel("Temperatura [K]")
    plt.ylabel("y [m]")
    plt.title(title_suffix)
    plt.grid(True)
    plt.legend()
    plt.tight_layout()
    plt.savefig(os.path.join(output_dir, filename), dpi=300)
    plt.close()


def plot_horizontal_profile(y_target, title_suffix, filename):
    plt.figure(figsize=(7, 5))
    for N in mesh_list:
        x = results[N]["x"]
        y = results[N]["y"]
        T = results[N]["T"]
        iy = np.argmin(np.abs(y - y_target))
        plt.plot(x, T[iy, :], label=f"{N}x{N}")
    plt.xlabel("x [m]")
    plt.ylabel("Temperatura [K]")
    plt.title(title_suffix)
    plt.grid(True)
    plt.legend()
    plt.tight_layout()
    plt.savefig(os.path.join(output_dir, filename), dpi=300)
    plt.close()


plot_vertical_profile(Lx / 2.0, "Perfil em t=4 h, x=Lx/2", "perfil_x_Lx2_t4h.png")
plot_vertical_profile(Lx / 4.0, "Perfil em t=4 h, x=Lx/4", "perfil_x_Lx4_t4h.png")
plot_horizontal_profile(Ly / 2.0, "Perfil em t=4 h, y=Ly/2", "perfil_y_Ly2_t4h.png")
plot_horizontal_profile(3.0 * Ly / 4.0, "Perfil em t=4 h, y=3Ly/4", "perfil_y_3Ly4_t4h.png")

# ============================================================
# EVOLUÇÃO TEMPORAL NOS PONTOS PEDIDOS
# ============================================================
plt.figure(figsize=(7, 5))
for N in mesh_list:
    plt.plot(results[N]["time"], results[N]["T_mid_hist"], label=f"{N}x{N}")
plt.xlabel("Tempo [s]")
plt.ylabel("Temperatura [K]")
plt.title("Evolução temporal em x=Lx/2, y=Ly/2")
plt.grid(True)
plt.legend()
plt.tight_layout()
plt.savefig(os.path.join(output_dir, "evolucao_temporal_centro.png"), dpi=300)
plt.close()

plt.figure(figsize=(7, 5))
for N in mesh_list:
    plt.plot(results[N]["time"], results[N]["T_q_hist"], label=f"{N}x{N}")
plt.xlabel("Tempo [s]")
plt.ylabel("Temperatura [K]")
plt.title("Evolução temporal em x=Lx/4, y=Ly/4")
plt.grid(True)
plt.legend()
plt.tight_layout()
plt.savefig(os.path.join(output_dir, "evolucao_temporal_quarto.png"), dpi=300)
plt.close()

print("\nSimulação concluída.")
print(f"Os gráficos foram salvos na pasta: {output_dir}")
\end{lstlisting}

\section{Comentários finais sobre a implementação}
\begin{itemize}[leftmargin=1.2cm]
    \item O código foi estruturado para resolver automaticamente as quatro malhas pedidas: $10\times 10$, $20\times 20$, $40\times 40$ e $80\times 80$.
    \item O passo de tempo é calculado a partir do limite de estabilidade do esquema explícito em 2D.
    \item A fronteira com radiação é tratada por linearização do termo $T^4$, o que torna a implementação mais simples e numericamente estável.
    \item Todos os gráficos são salvos em formato PNG, facilitando a inserção posterior no relatório final em PDF.
\end{itemize}

\section{Arquivos gerados pelo código}
Ao executar o script em Python, serão gerados arquivos PNG para:
\begin{itemize}[leftmargin=1.2cm]
    \item campos de temperatura em $t=0$, $1800$, $3600$, $7200$, $10800$ e $14400\,\si{s}$ para cada malha;
    \item perfis de temperatura em $t=4\,\si{h}$ para as posições pedidas;
    \item evolução temporal da temperatura em $(L_x/2, L_y/2)$;
    \item evolução temporal da temperatura em $(L_x/4, L_y/4)$.
\end{itemize}

\section{Conclusão}
A formulação apresentada atende aos itens solicitados no enunciado: equação governante com propriedades constantes e variáveis, discretização por diferenças finitas de segunda ordem, algoritmo explícito no tempo e implementação computacional em Python. A estrutura do código permite reproduzir diretamente os resultados gráficos exigidos na avaliação.

\end{document}
